\section{Introduction}

Courts evaluating gerrymandering claims need p-values, not percentile positions.
``More compact than 99.7\% of valid plans'' is not a standard statistical statement ---
it does not account for the precision of the 99.7\% estimate, which depends on
the ensemble's effective sample size.

The G.4 analysis established that a 1{,}000-step GerryChain ensemble for NC has
$\ess \approx 70$ (accounting for lag-1 autocorrelation $\rho_1 \approx 0.87$).
This means the ensemble provides approximately 70 independent plan observations,
not 1{,}000.
The nominal percentile of $\hat{p} = 0.003$ (3 plans more compact than the enacted
plan out of 1{,}000) has substantial uncertainty when derived from only 70
effective observations.

The permutation test framework corrects this: it converts the raw percentile into
an ESS-corrected p-value using the Beta posterior introduced in S.2.
The implementation is in \texttt{bisect-analysis::permutation::permutation\_test\_lower\_tail},
which takes a query plan's test statistic, the ensemble plans, and the ESS estimate
from G.4, and returns a \texttt{PermutationTestResult} with both raw and corrected p-values.

\paragraph{Null hypothesis.}
$H_0$: the query plan's test statistic is drawn from the same distribution
as the ensemble plans (the query plan is a random draw from the plan space).

\paragraph{Test statistic.}
We use the normalised edge-cut fraction $\mathcal{E}$ (lower = more compact)
as the primary test statistic, consistent with G.1.
